% Direct-authority diplomatic French for printed p.136.
\SourcePage{136}

\LayerTarget{tohoku:thm:1.10.1:p2}
\noindent de $\mathbf{C}$ dans $\mathbf{C}$, et un homomorphisme $f$ du foncteur identique dans
$M$, tels que pour tout $A\in \mathbf{C}$, $M(A)$ soit injectif et $f(A)$ soit un
monomorphisme de $A$ dans $M(A)$. La démonstration étant essentiellement
connue, nous esquisserons seulement les points principaux.

\LayerTarget{tohoku:lemma:1.10.1:1}
\textsc{Lemme 1.} \textit{Si $\mathbf{C}$ satisfait l'axiome \textup{AB 5)}, alors
l'objet $M\in \mathbf{C}$ est injectif si et seulement si pour tout sous-truc $V$ du
générateur $U$, et tout morphisme $v$ de $V$ dans $M$, $v$ se prolonge en un
morphisme de $U$ dans $M$.}

\LayerTarget{tohoku:lemma:1.10.1:1:proof}
Il suffit de prouver que la condition est suffisante, soit donc alors $u$ un
morphisme d'un sous-truc $B$ d'un $A\in \mathbf{C}$ dans $M$, montrons qu'il existe un
morphisme de $A$ dans $M$ prolongeant $u$. Considérons l'ensemble $P$ des
prolongements de $u$ à des sous-trucs de $A$ contenant $B$ (c'est bien un
ensemble, car en vertu de l'existence d'un générateur, les sous-trucs de $A$
forment un ensemble). Ordonnons-le par la relation de prolongement. En vertu
de la deuxième formulation de l'axiome AB 5) (cf. 1.5) cet ensemble est
inductif. Il admet donc un élément maximal, on est donc ramené au cas où $u$
est lui-même maximal, et à prouver qu'alors $B=A$.
Prouvons donc que si $B\neq A$, il existe un prolongement de $u$ à un
$B'\neq B$. Soit en effet $j$ un morphisme de $U$ dans $A$ tel que
$j(U)\not\subset B$, posons $B'=j(U)+B$ (donc $B'\neq B$). Soit
$V=j^{-1}(B)$ l'image inverse de $B$ par $j$, soit
$j':V\to B$ le morphisme induit par $j$, considérons la suite de
morphismes
$V\mathrel{\mathop{\to}^{\varphi'}}U\times B  \mathrel{\mathop{\to}^{\varphi}}B'\to0,$ oû le morphisme $\varphi'$ a pour composantes l'application
identique de $V$ dans $U$ et $-j'$, et $\varphi$ a pour composantes $j$ et
l'application identique de $B$ dans $B'$. On voit aussitôt que cette suite est
exacte, donc pour définir un morphisme $v$ de $B'$ dans $M$, il suffit de
définir un morphisme $w$ de $U\times B$ dans $M$ tel que
$w\varphi'=0$. Or soit $k$ une extension de $uj':V\to M$ à $U$,
et prenons pour $w$ le morphisme de $U\times B$ dans $M$ dont les composantes
sont $k$ et $u$, on vérifie aussitôt que $w\varphi'=0$, et que le morphisme
$v:B'\to M$ défini par $w$ prolonge $u$, ce qui achève la
démonstration du lemme 1.

\LayerTarget{tohoku:thm:1.10.1:construction:M1}
Soit $A\in \mathbf{C}$, soit $I(A)$ l'ensemble de tous les morphismes $u_i$ de
sous-trucs $V_i$ de $U$ dans $A$. Considérons le morphisme
$\oplus V_i\to A\times U^{(I(A))}$ dont la restriction au facteur $V_i$ a pour composantes
$-u_i:V_i\to A$, et l'injection canonique de $V_i$ dans le $i$.ème
facteur de la somme directe $U^{(I(A))}$. Soit $M_1(A)$ le conoyau du morphisme
envisagé, $f(A):A\to M_1(A)$ le morphisme induit par
l'épimorphisme canonique de $A\times U^{(I(A))}$ sur son quotient. Alors
$f(A)$ est un monomorphisme (vérification facile grâce au fait que le morphisme
canonique $\oplus V_i\to U^{(I(A))}$ est un monomorphisme en
vertu de AB 4)) et de plus tout morphisme $u_i:V_i\to A$ se
``prolonge'' en un morphisme $U\to M_1(A)$ (savoir le morphisme
induit sur le $i$.ème facteur de $U^{(I(A))}$ par l'épimorphisme canonique de
$A\times U^{(I(A))}$ sur son quotient $M_1(A)$).
\LayerTarget{tohoku:thm:1.10.1:construction:transfinite:p1}
Définissons par récurrence transfinie, pour tout nombre ordinal $i$, l'objet
$M_i(A)$, et pour deux nombres ordinaux $i\leqq j$, un morphisme
\textit{injectif} $M_i(A)\to M_j(A)$, de telle façon que les
$M_i(A)$ pour $i<i_0$ ($i_0$ nombre ordinal fixé) forment un système inductif.
Pour $i=0$, on prendra $M_0(A)=A$ ; pour $i=1$, $M_1(A)$ et
$M_0(A)\to M_1(A)$ sont déjà définis. Si la construction est
faite pour les ordinaux $<i$, et si $i$ est de la forme $j+1$,

\TohokuContinuation{137}
